% Generated by Listas/make_statements.py from pratica-lecture-05.md, item 14. Do not edit.
\textbf{Source.} Muñoz-Garcia \& Toro-Gonzalez (2019), \emph{Strategy and Game Theory: Practice Exercises with Answers}, Ch. 6, Exercise 6 ``Collusion in Donations in an Infinitely Repeated Public Good Game'', pp.~292--293 {[}PDF 303--304{]}. Answer in the source: yes, a worked answer follows the exercise. Reproduced verbatim (statement only; the book's worked answer follows it; ``Grim-Triger'' is the book's spelling).



\begin{quote}

Consider a public good game between two players, A and B. The utility function of every player \(i = \{A, B\}\) is



\[u_i(g_i) = (w - g_i)^{1/2}\left[m\left(g_i + g_j\right)\right]^{1/2}\]



where the first term, \(w - g_i\), represents the utility that the player receives from money (i.e., the amount of his wealth, \(w\), not contributed to the public good). The second term indicates the utility he obtains from aggregate contributions to the public good, \(g_i + g_j\), where \(m \geq 0\) denotes the return from aggregate contributions. For simplicity, assume that both players have equal wealth levels \(w\).



\begin{enumerate}

\def\labelenumi{(\alph{enumi})}

\item

  Find every player \(i\)'s best response function, \(g_i(g_j)\). Interpret.

\item

  Find the Nash Equilibrium of this game, \(g^*\).

\item

  Find the socially optimal contribution level, \(g^{SO}\), that is, the donation that each player should contribute to maximize their joint utilities. Compare it with the Nash equilibrium of the game. Which of them yields the highest utility for every player \(i\)?

\item

  Consider the infinitely repeated version of the game. Show under which values of players' discount factor \(\delta\) you can sustain cooperation, where cooperation is understood as contributing the socially optimal amount (the one that maximizes joint utilities) you found in part (c). For simplicity, you can consider a ``Grim-Triger strategy'' where every player starts cooperating (that is, donating \(g^{SO}\) to the public good), and continues to do so if both players cooperated in all previous periods. However, if one or both players deviate from \(g^{SO}\), he reverts to the Nash equilibrium of the unrepeated game, \(g^*\), thereafter.

\end{enumerate}

\end{quote}

